% vim: ai sw=2 spell spelllang=cs

\begin{frame}
  \frametitle{S čím budeme pracovat?}

  \heading{Software}
  \begin{itemize}

    \item GIT
    \begin{itemize}
      \item Úvod do GITu dnes
      \item Podrobněji: \url{http://book.git-scm.com/}
    \end{itemize}

    \pause

    \item Zdroje jádra
    \begin{itemize}
      \item Použijeme předinstalované z RPM (\file{/usr/src/kernels/})
      \item \url{http://git.kernel.org/} (od 2.6.12)
      \item \texttt{full-history-linux} GIT (od 0.01 do 2.6.26)
    \end{itemize}

    \pause

    \item \textsc{Qemu}
    \begin{itemize}
      \item HW k práci s PCI, I/O, na přerušení, DMA\dots
    \end{itemize}
  \end{itemize}

\end{frame}

\section{GIT}
\frame{\sectionpage}

\begin{frame}
  \frametitle{Úkol}

  \heading{GIT a repozitář PB173}
  \begin{enumerate}
    \iffimuni {
      \item Stáhněte si repozitář na aisu (doporučeno, můžete vynechat)
      \begin{itemize}
	\item \cmd{git clone --bare} \url{git://github.com/jirislaby/pb173}
      \end{itemize}

      \item Vytvořte si lokální klon
      \begin{itemize}
	\item \cmd{git clone xlogin@aisa:pb173}
      \end{itemize}
    } \else {
      \item Stáhněte si repozitář
      \begin{itemize}
	\item \cmd{git clone} \url{git://github.com/jirislaby/pb173}
      \end{itemize}
    } \fi

    \item Prozkoumejte strukturu
    \begin{itemize}
      \iffimuni
	\item Příklady ze cvičení
	\item Adresář pro domácí úkoly
      \else
	\item Příklady ke kurzu
      \fi
    \end{itemize}

  \end{enumerate}
\end{frame}

\begin{frame}[fragile]
  \frametitle{Úkol}

  \heading{GIT a úpravy souborů}
  \begin{enumerate}
    \item Změňte cokoliv v souboru \cmd{sandbox/hello}
    \item Zkontrolujte změny (\cmd{git diff --color})

    \item Uložte do lokálního repozitáře (\cmd{git commit -a})
    \begin{itemize}
      \item Git může chtít nastavit jméno a e-mail (instrukce jsou na stdout)
      \item Formát logu (\emph{vzor:} \href{https://git.kernel.org/cgit/linux/kernel/git/torvalds/linux.git/log}{\nolinkurl{git.kernel.org}})
      \begin{semiverbatim}
	Shrnutí na řádek
	<Volný řádek>
	Odůvodnění (delší popis)
	<Volný řádek>
	Podpisy a CC
      \end{semiverbatim}\vspace{-\baselineskip}
    \end{itemize}

    \pause

    \item Smažte \cmd{sandbox/hello} (\cmd{git rm})
    \item \cmd{git commit -a}

  \end{enumerate}

\end{frame}

\begin{frame}
  \frametitle{GIT a commity}

  \begin{itemize}
    \item Každý commit je pojmenován SHA hashem
    \begin{itemize}
      \item Např.\ \code{d3323c1503d54d83b0eae6c7927dede2d2973059}, lze i
	zkráceně \code{d3323c1503}
    \end{itemize}

    \pause

    \item \cmd{HEAD} (\cmd{@}) je alias pro horní/poslední commit

    \pause

    \item Lze odkazovat předchůdce pomocí \textasciitilde\ nebo \textasciicircum
    \begin{itemize}
      \item Např.\ \code{HEAD}\textasciicircum\textasciicircum\textasciicircum,
	\code{HEAD~1}, \code{d3323c1503~5}
    \end{itemize}

    \pause

    \item Lze vytvořit jmenný alias, tzv.\ tag

    \pause

    \item Více: \doc{man gitrevisions}

  \end{itemize}

\end{frame}

\begin{frame}
  \frametitle{Úkol}

  \heading{GIT a commity}
  \begin{enumerate}
    \item Zkontrolujte log, zda obsahuje 2 změny (\cmd{git log --color})
    \item Podívejte se na poslední 2 změny (\cmd{git show --color HEAD},
      resp.\ \cmd{HEAD^1})

    \pause

    \item Vygenerujte záplaty ze 2 posledních commitů (\cmd{git format-patch
      -2})

    \iffimuni
      \pause

      \item Proveďte push (\cmd{git push}, jen pokud máte klon na aise)
    \fi

  \end{enumerate}

  \iffimuni {
    \bigskip
    \pause

    \begin{center}
      \heading{Odevzdávání domácích úkolů odesláním záplaty na e-mail. \\
	Posílejte jen plain-text přílohy!}
    \end{center}
  } \fi

\end{frame}

\begin{frame}
  \frametitle{Záplatování}

  \heading{Záplaty v linuxovém jádře}
  \begin{enumerate}
    \item Poslat odpovídajícímu správci
    \begin{itemize}
      \item \file{scripts/get_maintainer.pl}
      \item \cmd{git send-email}
    \end{itemize}

    \pause

    \item Poslat na ML ,,pull request'' a vystavit celý GIT strom
    \begin{itemize}
      \item \cmd{git push}
      \item \cmd{git request-pull}
    \end{itemize}
  \end{enumerate}

  \pause

  \begin{itemize}
    \item Zevrubný popis v \doc{Documentation/SubmittingPatches}
  \end{itemize}

\end{frame}


\section{Jádro}
\frame{\sectionpage}

\begin{frame}
  \frametitle{Jádro vs.\ programy}

  \heading{Hlavní rozdíly}
  \begin{itemize}
    \item Žádné knihovny
    \begin{itemize}
      \item Zejména \emph{libc} (\code{printf}, \code{strlen}, \code{malloc},
	\dots)
      \item Ale ani ostatní (např.\ \file{m}, \file{pthread}, \file{ssl},
	\file{z})
    \end{itemize}

    \pause

    \item Ne/oddělený \emph{paměťový prostor}
    \begin{itemize}
      \item Chybový kód může přepisovat cokoliv (i FW)
    \end{itemize}

    \pause

    \item \emph{Počáteční funkce} (\code{main})
    \begin{itemize}
      \item \code{module\_init} (=\code{main}), \code{module_exit}
	(=\code{on_exit})
    \end{itemize}

    \pause

    \item Pád jádra $\Rightarrow$ pád všeho

  \end{itemize}
\end{frame}

\begin{frame}
  \frametitle{Kód}

  \begin{itemize}
    \item Výhradně \emph{GNU C} (\xes už jen \textsc{gcc} $\geq$ 4.x)

    \pause

    \item Pevně daný \emph{CodingStyle}
    \begin{itemize}
      \item Kontrola: \cmd{scripts/checkpatch.pl} (není 100\%)
    \end{itemize}

    \pause

    \item Vše slinkováno do jednoho celku
    \begin{itemize}
      \item \file{vmlinux} $\rightarrow$ \file{(b)zImage}
      \item Ale moduly (standardní ELF: \file{*.ko})
    \end{itemize}
  \end{itemize}

\end{frame}

\begin{frame}
  \frametitle{Úkol}

  \heading{Průzkum modulu (\file{.ko} objektu)}
  \begin{enumerate}
  \item Zvolte si jeden modul v systému (\cmd{lsmod})
  \item Zavolejte na něm \cmd{modinfo}
  \item Proveďte \cmd{objdump -d} sekce \file{.modinfo}
  \item Porovnejte oba výstupy
  \end{enumerate}
\end{frame}

\begin{frame}
  \frametitle{Dokumentace}

  \begin{itemize}
    \item Linux Device Drivers, 3. edice
    \begin{itemize}
      \item Ke stažení: \url{https://lwn.net/Kernel/LDD3/}
    \end{itemize}

    \item \doc{Documentation/*}
    \item Generovaná tamtéž (podobná \textsc{DocBook})
    \item Kód
    \begin{itemize}
      \item Hledání: \url{http://lxr.free-electrons.com/ident}
    \end{itemize}
  \end{itemize}

  \medskip

  \heading{Demo:} pomocí lxr najděte v jádře obecnou (v \file{lib/}) a
  optimalizovanou pro \xes 32-bity (v \file{arch/x86/lib/}) implementaci
  \code{strlen}
\end{frame}

\begin{frame}
  \frametitle{Úkol}

  \heading{Hello World}
  \begin{enumerate}
    \item Spusťte virtuální stroj

    \item Prozkoumejte adresář \file{01} z pb173 git repozitáře
    \begin{itemize}
      \item \file{Makefile}, \file{pb173.c}
    \end{itemize}

    \item Do \code{init} funkce doplňte výpis ,,Hello World''
    \begin{itemize}
      \item \code{printk(KERN_INFO "...\\n");} \emph{(ne, není tam čárka)}
      \item Nebo zkráceně: \code{pr_info("...\\n");}
      \item Objeví se v \file{/proc/kmsg} a na konzoli
    \end{itemize}

    \item Přeložte a vložte do systému
    \begin{itemize}
      \item \cmd{make}, \cmd{insmod <modul.ko>}
    \end{itemize}

    \item Zkontrolujte výstup \cmd{dmesg}
    \item Můžete udělat \cmd{commit} \iffimuni a \cmd{push} \fi do svého
      repozitáře
  \end{enumerate}
\end{frame}

\iffimuni
\subsection{Paměť}
\fi

\begin{frame}
  \frametitle{Paměť}

  \begin{itemize}
  \item \pozor: omezený zásobník (4--8\,K)
  \begin{itemize}
  \item Žádná nebo malá rekurze
  \item Jen pro malá data (\code{int}, malé \code{struct}, krátká pole, \dots)
  \end{itemize}

  \pause

  \item Pracuje se se stránkami (na \xes 4\,K--1\,G)
  \begin{itemize}
  \item Omezená velikost alokace (fragmentace)
  \item Podrobnosti \iffimuni v dalších cvičeních \else dále \fi
  \end{itemize}
  \end{itemize}
\end{frame}

\begin{frame}[fragile]
  \frametitle{Jednoduché alokace}

  \heading{Malé alokace (sta bajtů, maximálně až cca.\ 4\,M)}
  \begin{itemize}
  \item Horní mez závislá na architektuře
  \item \code{kmalloc}, \code{kfree} (\header{linux/slab.h})
  \end{itemize}

  \pause
  \bigskip

  Alokace mají (většinou) GFP parametr. Ten určuje, co si alokátor může dovolit
  (spát, swapovat, použít HIGHMEM, \dots). Prozatím nám stačí
  \code{GFP_KERNEL}.

  \pause
  \bigskip

  \begin{lstlisting}
void *mem = kmalloc(100, GFP_KERNEL);
if (mem) {
  ...
}
kfree(mem);
  \end{lstlisting}

\end{frame}

\begin{frame}
  \frametitle{Úkol}

  \begin{enumerate}
  \item V \code{exit} funkci naalokujte 1000 bytů (\code{kmalloc})
  \item Udělejte do nich \code{strcpy} ,,Bye''
  \item Vypište paměť jako řetězec (\code{printk} a \code{\%s})
  \item Uvolněte paměť
  \item Vložte a odeberte modul ze systému (\cmd{insmod} a \cmd{rmmod})
  \item \cmd{dmesg}
  \item Můžete provést \cmd{commit} \iffimuni a \cmd{push} \fi
  \end{enumerate}
\end{frame}

